\chapter{Discussion} \label{chap:discussion}

This chapter interprets the results of Chapter~\ref{chap:results} hypothesis by hypothesis, explains why each part of EADD behaves as it does, places the detector among the methods reviewed in Chapter~\ref{chap:literature}, and states the limitations of the study. All numbers refer to the final registered test (seeds 3000 to 3009) unless stated otherwise.

\section{H1: Why a Reservoir Reference Finds Slow Drift}
\label{sec:disc-h1}

On the 14 streams with slow feature drift, EADD detected every drift point (recall 1.00), while D3-LightGBM detected 0.53 and D3 0.48, and EADD's mean time to detection (MTD) was 443 samples against 879 and 1,008. The gap was largest on gradual drift (recall 1.00 against 0.20 for D3-LightGBM, MTD 582 against 1,373) and smaller on incremental drift (1.00 against 0.85, MTD 304 against 385). The part-by-part benchmark located the source of this gain in the reference window: switching on the reservoir alone in EADD's own code, which otherwise reproduces D3-LightGBM, raised the gradual recall from 0.00 to 0.50 and the incremental recall from 0.88 to 1.00.

The explanation lies in what each reference contains while a slow drift is under way. A sliding reference holds the 500 samples just before the current window, so it moves with the drift. When the share of new samples rises over 1,000 samples, or the mean moves linearly over 1,000 samples, the reference and the current window are only a few hundred samples apart and differ by a small fraction of the full shift at every check. A classifier two-sample test (C2ST) \citep{LopezPazOquab2017} on such windows rarely reaches an AUC of 0.7, and by the time the drift ends the reference has already absorbed the new distribution. A reservoir reference is a uniform sample of all data since the last alarm \citep{Vitter1985}, so during a slow drift it is still dominated by the old distribution, and the contrast with the current window grows towards the full shift. The detector therefore sees the accumulated change rather than the change within one window length.

On abrupt drift the two references contain the same thing at the moment of change, namely data from before the drift, and the current window fills with new data at the same rate. This is why EADD gained nothing there but lost nothing either: its abrupt MTD was 115 samples against 122 for D3-LightGBM, within the registered margin of 50 samples (Holm $p = 0.023$). The result supports the revised H1 for slow drift, and it holds against the original D3 with logistic regression as well.

The same mechanism has a cost. When both detectors found an incremental drift, D3-LightGBM found it sooner (a delay of 193 against 304 samples over detected drift points only). Two explanations are consistent with this, and the experiments did not separate them. First, the delay over detected points is computed on different sets: D3-LightGBM missed 15\% of the incremental drift points, and the points it found are likely the easier ones, while EADD's mean includes the harder points that only it detected. Second, an EADD alarm needs the permutation test to reject as well as the AUC to pass the threshold, and once a slow drift is under way, the most recent windows in the history null are themselves partly drifted, which widens the null and can delay the rejection. The practical reading is that the reservoir makes slow drift detectable at all, while a sliding reference, when it detects such a drift, can be faster.

\section{H2: Why the History Null Cuts False Alarms}
\label{sec:disc-h2}

On the six autocorrelated streams without drift, EADD raised 5.59 false alarms per 10,000 samples against 19.5 for D3-LightGBM and 18.5 for D3, fewer on every stream (Holm $p = 0.023$ against D3-LightGBM, the primary comparison, and 0.016 against D3). On the other streams without drift, EADD raised 0.75 against 2.44. Since every alarm on these streams is false, they measure false alarms directly, which supports H2. Addendum 7b isolates the cause: the reservoir with the label-shuffle null raised as many false alarms as D3-LightGBM (19.5 per 10,000 samples), while a sliding reference with the history null raised 4.15, so the history null, not the reservoir, removes them.

The reason is that two windows of an autocorrelated stream differ by chance far more than two random halves of the same data. For an AR(1) process with $\rho = 0.9$, the variance of the mean of 200 samples is about $\frac{1}{n}\frac{1+\rho}{1-\rho}$, a standard deviation of about 0.31 against 0.07 for independent samples. A sliding-reference detector without a test, such as D3, alarms whenever the stream wanders this far, and a permutation test that shuffles the window labels assumes exchangeable samples \citep{Good2005}, so its null is as narrow as for independent data and it rejected every test at $\rho = 0.9$ in Chapter~\ref{chap:methodology}. The history null instead compares earlier 200-sample windows of the same stream with the reference, so its null AUCs already contain the wandering that autocorrelation causes. An observed AUC is then judged against what this stream produces without drift, and the threshold of 0.7 is no longer the only barrier against chance.

The guarantee of this test is limited: it controls the false-rejection rate of each single test: in 1,000 tests per condition on no-drift data, 3.8\% to 6.9\% rejected at $p \le 0.05$, and 2.9\% to 4.2\% at EADD's alarm rule in steady state. It does not control the error over a whole stream. EADD tests every 50 samples, so false alarms still accumulate with the length of the stream, and the 5.59 false alarms per 10,000 samples on autocorrelated streams are the expected result of many tests at a fixed per-test level rather than a failure of the test. The history null also assumes that the dependence in the stream's recent past is representative of the present, which is weakest right after an alarm, when few earlier windows are available (Limitation~1 in Section~\ref{sec:limitations}).

\section{H3: Why ASPT Keeps Detection Quality}
\label{sec:disc-h3}

The adaptive sequential permutation test (ASPT) gave the same recall (1.00 against 1.00), MTD (333.7 against 333.5) and false alarms (0.98 against 0.96 per drift stream, 0.75 against 0.63 per 10,000 samples on no-drift streams) as the full permutation test, all equivalent within the registered margins ($p < 0.01$), and it used 52 instead of 163 permutations per test. On the INSECTS streams it cut the runtime of a run from 453 to 210 seconds (54\% less) with identical detections in 20 of 20 run pairs. H3 is supported.

The saving follows from what the decision needs. The detector only needs to know whether the $p$-value is below $\alpha$, not its exact value. When the observed AUC lies far in the tail, as it does for most AUCs above 0.7, no null draw reaches it, and after 59 draws the bound $1/(i+1) \le \alpha/3$ already allows a rejection. When the observed AUC lies inside the null, more than $2\alpha$ of the draws soon reach it, and the test accepts early. Only borderline cases need all 199 draws. This is the idea of sequential Monte Carlo tests \citep{BesagClifford1991, Gandy2009}, and splitting $\alpha$ between the early and the final decision keeps the simulated Type~I error below the nominal level (0.031, where the earlier rule had 0.059). Because borderline cases are rare at a threshold of 0.7, the early decisions almost never differ from the full test, which is why the detections were the same.

\section{H4: SHAP Attribution in a Classifier-Based Detector}
\label{sec:disc-h4}

At the first detection of each drift point, the TreeSHAP ranking \citep{Lundberg2017, LundbergInteraction2020} placed the truly drifting features above chance: precision@$k$ was 0.75 against 0.31 by chance, and the top-ranked feature was a drifting one in 98.4\% of detections ($p < 0.001$). Without PokerHand and Rialto, where some of the shifted features are constant before the drift (which makes those drifts easy to detect), precision@$k$ was 0.80 against 0.32. H4 is supported.

The result is expected for any classifier-based detector, not only for EADD. A classifier trained to separate the reference from the current window can only use features whose distribution differs between them, so an attribution of its predictions should concentrate on those features. Any classifier-based detector, D3 included, could attach TreeSHAP or another attribution method to its discriminator. This thesis shows empirically, with EADD as the detector, that such an attribution does name the drifting features on streams where they are known; whether the size of the effect is the same with a logistic regression discriminator was not measured. The ranking says which features separate the windows, not why they changed, so it is a lead for investigation rather than a causal diagnosis \citep{CausalDrift2025}, in line with other work that explains drift with Shapley values \citep{DBShap2024, HinderFeature2024}.

The comparison with univariate rankings (Section~\ref{sec:exp4}) limits what H4 means in practice. On these streams a per-feature KS statistic between the same two windows named the drifting features more accurately than TreeSHAP (precision@$k$ 0.97 against 0.75). TreeSHAP ranks the features by their use in a classifier that also fits noise and correlated features, while the injected drifts change single marginal distributions, the case a univariate statistic is built for. TreeSHAP attribution is therefore above chance but not the best ranking for marginal drift; a monitor could report the KS ranking beside it, and the case of drift in the dependence between features remains to be tested.

\section{Position among Existing Detectors}
\label{sec:disc-position}

Against D3 \citep{Gozuacik2019}, the comparison is controlled: EADD's own code with its parts switched off reproduces D3-LightGBM (recall 0.84 against 0.84, MTD 299 against 302 on 20 benchmark streams), so the differences reported here come from the reservoir reference and the history-null test, not from the implementation. The classifier matters too: replacing logistic regression with LightGBM raised the incremental recall of D3 from 0.72 to 0.88 in the benchmark. In the benchmark of \citet{Lukats2024}, D3 was among the strongest fully unsupervised detectors, with the best mean time ratio on INSECTS Abrupt. The present results agree that a sliding-reference C2ST is hard to beat on abrupt drift and show where it falls short, namely on slow drift and on autocorrelated streams.

The newer online detectors MMD-Online, LSDD-Online \citep{CobbSequential2022} and CVM-Online \citep{RossAdams2012}, run from Alibi Detect \citep{AlibiDetect} and calibrated to an expected run time of 2,500 samples, show a different trade-off. EADD's own figures in these runs (seeds 3000 to 3002) differ slightly from those of the main test. The online detectors were faster: on slow drift CVM-Online had an MTD of 224 samples and LSDD-Online 365 against EADD's 443, and on abrupt drift LSDD-Online had 66 and CVM-Online 78 against 114. They also raised far more false alarms: on slow-drift streams 6.88 to 9.88 per stream against 1.52 for EADD (about 4.5 to 6.5 times as many), 6.82 to 7.13 per 10,000 samples on no-drift streams against 0.75, and 18.0 to 19.2 on autocorrelated streams against 5.89. In the Friedman ranks over the 21 drift streams, EADD had the best mean recall rank (2.83), CVM-Online the best MTD rank (1.90), and EADD's false-alarm rank was better than that of each online detector by more than the critical difference. An expected run time of 2,500 samples implies a few false detections per 10,000 samples by design, so this trade-off depends on the calibration, and the pairwise tests were not significant with three seeds (smallest Holm $p = 0.19$). EADD is therefore not the fastest detector, but at this calibration of the online detectors it combined a high recall with few false alarms.

\section{What the Recent Literature Means for These Results}
\label{sec:disc-literature}

The clearest lesson from work published in 2025 and 2026 is the difference between per-test and stream-wide error control. ASPT is sequential over the permutations of one window, not over time, so EADD's level holds for each test. Betting-based sequential two-sample tests \citep{PodkopaevRamdas2023} and e-values for the C2ST \citep{EC2ST2024} control the probability of any false alarm at any stopping time, and they also use a classifier as the test statistic, so they are a natural route to giving a detector like EADD stream-wide guarantees.

On evaluation, the framework of \citet{CerqueiraKDD2026} recommends injecting controlled drifts into real data, measuring timing-aware metrics and keeping tuning out of the test data. The design of this thesis follows the same principles: drifts injected into six real bases with known drift points and drifting features, detection windows that count a late detection as a miss, and fixed settings for all detectors. Tuning was not fully separated from the test data: EADD's history stride and alpha split were chosen on benchmark streams built from the same real bases.

Several recent detectors address settings outside this thesis. PUDD \citep{PUDD2025} monitors the prediction uncertainty of a deployed model, which EADD does not need. DriftLens \citep{DriftLens2024} works on deep-learning embeddings of unstructured data. The Gaussian Split Detector \citep{GaussianSplit2024} and related detectors try to separate real drift, which moves the decision boundary, from virtual drift, while EADD by design reports any change in $P(X)$. The autoencoder detector of \citet{Komorniczak2026} targets tabular streams in which new classes appear but has no released code. These methods complement rather than replace a C2ST on the input features.

\section{Threats to Validity Addressed by Design}
\label{sec:disc-validity}

Several threats that weakened the earlier evaluation were removed by the design of the final test. All detectors used identical settings: windows of 500 and 200 samples, a threshold of 0.7, 5-fold cross-validation, the same LightGBM parameters, a check every 50 samples, the same stream objects, seeds and scoring code, and a random start offset between 0 and 49 that was the same for all detectors, so no detector benefited from the check grid. The unit of analysis was the stream rather than the run, which avoids treating seeds of the same stream as independent evidence, and the tests used Holm correction within each registered family.

Each test's hypotheses, data, seeds and decision rule were committed to the code repository before its runs, and each test used fresh seeds (1000 to 1009 for the confirmatory test, 3000 to 3009 for the final test), so no result was chosen among several runs. H1, H2 and the detection part of H3 were re-specified after earlier registered versions failed, and each new version was tested on new seeds. The four flaws of the earlier detector were fixed as options and unit-tested, while the original thesis runs still reproduce with the old settings (Chapter~\ref{chap:methodology}).

Three reviews, each working from the raw run files and carried out with the help of AI assistants (Section~\ref{sec:implementation}), checked the evaluation before the final test. The code review confirmed that the four fixes are correct and that no information leaks between the windows, and it found that D3 and D3-LightGBM checked every 51 samples instead of 50; this was fixed by letting every detector check every 50 samples from a random start per seed. The statistics review recomputed every registered number with independent code, and all matched. The data review confirmed the drift patterns, the drift points and the ground truth of the drifting features, and that no labels reach any detector; it found the shifted features that are constant before the drift on PokerHand and Rialto, so the results are reported with and without these two bases. The code, the raw runs and the registration documents are public, so every number in Chapter~\ref{chap:results} can be recomputed.

\section{Limitations}
\label{sec:limitations}

\textbf{Limitation 1: level of the test.} The history-null test is close to its nominal level but not exactly at it. At $\alpha = 0.05$ it rejected 3.8\% to 6.9\% of no-drift tests, and three of ten conditions exceeded the bound of 0.075: AR 0.3 and AR 0.5 in the steady state and shuffled Electricity right after an alarm. Right after an alarm only 31 earlier windows exist, so an alarm needs every null AUC to stay below the observed one, a level of 1/32; the after-alarm rates of 3.8\% to 6.9\% are above that. At AR 0.9 the confirmatory test measured 7.0\% in the steady state and 9.4\% after an alarm. Two fixes tried on development data did not help. In use this means more false alarms on strongly autocorrelated streams and soon after an alarm than the nominal level suggests; a longer wait before the first test after an alarm or a betting-based test could address it.

\textbf{Limitation 2: scope of the advantage.} On abrupt drift EADD is equivalent to D3-LightGBM, not better, and when both detect an incremental drift, D3-LightGBM is faster (193 against 304 samples delay when detected). On the four annotated INSECTS streams of the benchmark (Section~\ref{sec:insects}), EADD was not better: its recall was lower on the two incremental streams (0.80 and 0.60 against 1.00) and it raised more alarms counted as false (18.8 and 20.2 against 13.4 and 14.2 per stream). The INSECTS data change slowly between the documented drift points, and the reservoir's sensitivity to slow change produces alarms there that the annotation counts as false. EADD is therefore most useful where slow drift matters; streams with richer annotation would show how many of these alarms are real changes.

\textbf{Limitation 3: evidence base.} Fresh seeds are new random draws from the same generators and datasets, not unseen data sources, so the results generalise to new runs of these streams more firmly than to new domains. The newer online detectors were run on three seeds and calibrated at one expected run time; more seeds and several calibrations would give a fuller comparison.

\textbf{Limitation 4: scoring the attribution.} The accuracy of the feature attribution can be scored only where the drifting features are known, that is on the synthetic and injected streams. On real streams the ranking remains a lead; annotated real streams with documented changed features would allow it to be scored there as well.
